\section{A seed-to-refinement construction}\label{sec:cones}

The initial table has a large approximation factor.  We now describe one
refinement that reduces this factor while controlling expected excess over
the true cost.  Throughout the construction, fix the entire seed table;
randomness introduced by the refinement is independent of it.

\begin{theorem}[Refinement]\label{thm:refinement}
Use the parameters and interval tree of Section~\ref{sec:geometry} in the
large-input branch.  Let $A$ be a power of two with
$1\le A\le A_{\rm init}$, and let $U$ be an integer table on every labeled
tree node and every valid target state, with
\[
 0\le U(I,q)\le N,\qquad c(I,q)\le U(I,q)\le A c(I,q).
\]
Set $a=\max\{1,A/F\}$.  There is a randomized, finitely executable
refinement table $d_S$ such that, for every state $(I,q)$,
\begin{equation}\label{eq:refinement-guarantee}
 d_S(I,q)\ge c(I,q)/a,\qquad
 \E\bigl[\pos{d_S(I,q)-c(I,q)}\bigr]\le c(I,q)/\ell.
\end{equation}
The first inequality holds for every outcome of the refinement randomness;
the expectation in the second is conditional on the fixed table $U$.
At zero-cost states the returned value is zero for every outcome.
\end{theorem}

We call a seed \emph{good} if its input bounds
$c(I,q)\le U(I,q)\le A c(I,q)$ hold simultaneously at every state.
This section and the next define the refinement table
and establish its elementary
invariants.  Sections~\ref{sec:comparison}--\ref{sec:accuracy} prove the
expected-error bound.  Sections~\ref{sec:numerics} and~\ref{sec:work} give
the numerical and evaluation algorithms, including their work on
inaccurate seed tables.  The distinction matters: the refinement can
underestimate true cost when $a>1$.

\subsection{Cost cones and their local evaluation}

The number $a$ is an integer power of two.  Define
\[
 \widetilde U(I,q)=\max\{\ceil{U(I,q)/a},\abs{z(I,q)}\}.
\]
Since $Fc(I,q)$ is an integer, rounding does not spoil the bounds
\begin{equation}\label{eq:scaled-seed}
 \max\{c(I,q)/a,\abs{z(I,q)}\}
 \le\widetilde U(I,q)\le Fc(I,q).
\end{equation}
These bounds do not give endpoint regularity: neighboring seed entries
can have quite different overestimates. We replace the scaled table by
a lower envelope of cones. The charge $L_0|q-r|_1$ for moving away from
a center $r$ preserves the lower bounds $c/a$ and $|z|$, while making
the envelope Lipschitz. To evaluate it with few seed queries, we retain
centers on grids at several cost scales. A grid at the scale of the true
cost will provide a nearby inexpensive center; its value will then
force every minimizing center to be nearby as well.

The scales $b$ are the powers of two from $1$ to the largest power of two
at most $N$.  At scale $b$, put $g_b=\max\{1,b/P\}$.  A grid state has
both coordinates in $g_b\Z\cap[0,|y|]$, in nondecreasing order.
The state $r_b(q)$ rounds both coordinates of $q$ down to this grid.
All these spacings are integers.  If $g_b=1$, rounding is exact;
otherwise
\begin{equation}\label{eq:grid-error}
 \abs{q-r_b(q)}_1<2b/P.
\end{equation}

At a nonleaf $I$, consider all cones
\[
 q\longmapsto \widetilde U(I,r)+L_0\abs{q-r}_1
\]
whose centers are grid states at some scale $b$ with
$b/2\le U(I,r)\le4Ab$.  Include also the cone at every state $r$ with
$U(I,r)=0$, whether or not it occurs on one of these grids.  Define
$w_0(I,q)$ to be their minimum.  At a leaf set $w_0(I,q)=c(I,q)$.
The exact leaf rule will be used for every later table as well.

\begin{lemma}[Initial cones]\label{lem:initial-cones}
For the fixed seed in Theorem~\ref{thm:refinement},
\begin{equation}\label{eq:initial-cones}
 \max\{c/a,|z|\}\le w_0\le2Fc,
 \qquad \Lip(w_0(I,\cdot))\le L_0.
\end{equation}
At a nonleaf, the following local rule computes this global minimum
exactly.  If $U(I,q)=0$, return zero.  Otherwise omit the zero cones,
enumerate scales
\begin{equation}\label{eq:initial-local-scales}
 U(I,q)/(16A^2)\le b\le4A U(I,q),
\end{equation}
and enumerate grid centers at sup distance at most $U(I,q)$ from $q$,
retaining the eligibility test $b/2\le U(I,r)\le4Ab$.
\end{lemma}

\begin{proof}
Each cone dominates $c(q)/a$ and $|z(q)|$ by
\eqref{eq:scaled-seed} and the endpoint Lipschitz bounds.  Each has
Lipschitz constant $L_0$, so their minimum does too.  Leaves satisfy
these claims directly.

If $c(I,q)=0$, the seed has $U(I,q)=0$, and its own cone gives zero.
For $c=c(I,q)>0$, choose $b\le c<2b$ and set $r=r_b(q)$,
$e=|q-r|_1$.  By~\eqref{eq:grid-error},
\[
 c(I,r)\ge b-2b/P\ge b/2,
 \qquad c(I,r)<2b+2b/P\le4b.
\]
Thus $r$ is an eligible center.  Its cone has value at most
\[
 Fc(I,r)+L_0e
 \le Fc+(F+L_0)e
 \le Fc+66Fb/P\le2Fc.
\]

Any minimizing cone at a positive-cost state consequently has center
distance at most $2Fc/L_0=c/16$.  Its center's true cost lies in
$[15c/16,17c/16]$.  A zero center cannot minimize, since its distance
would be at least $c$.  For a minimizing nonzero center, eligibility
therefore gives
\[
 b\ge \frac{U(I,r)}{4A}
   \ge\frac{15c}{64A}\ge\frac{U(I,q)}{16A^2},
 \qquad
 b\le2U(I,r)\le\frac{17Ac}{8}\le4AU(I,q).
\]
Its sup distance is at most $c/16\le U(I,q)$.  Thus the local list
contains a global minimizer, and every member of the local list belongs
to the global family.  The two minima agree.
\end{proof}

We use this local rule as the executable definition even when the seed
is inaccurate, returning $N$ if its nonzero branch has no candidate.
Its work remains finite without any seed guarantee.  Indeed, the scale
interval has ratio at most $64A^3$, and $U(I,q)/b\le16A^2$ there.
The local list therefore contains
$O((1+PA^2)^2(1+\log A))=H^{O(p_*+1)}$ raw seed queries.
The returned values are nonnegative integers of magnitude
$O((L_0+1)N)$.  Global cone identities and Lipschitz continuity are
asserted only for a seed satisfying Theorem~\ref{thm:refinement}.

\subsection{Wide actions and Bellman propagation}

The next operator combines child estimates into a parent estimate. We
first allow each child to choose its target independently, paying the
connection charge for gaps and overlaps. Its minimum over earlier child
tables will remain an available parent option throughout the refinement;
the more restricted choices will use current child tables. Extending
each action cost by an endpoint cone retains spatial
regularity. Repeated propagation will also supply the average cost table
used to predict child alignments.

Let a nonleaf $I=(l,u]$ have child cuts
$l=l_0\le l_1\le\cdots\le l_M=u$, so
$I_i=(l_{i-1},l_i]$.  A \emph{representative} is a pair $(b,r)$ with
$r$ a grid state at scale $b$ and
\begin{equation}\label{eq:representative}
 b/(4a)\le w_0(I,r)\le8Fb.
\end{equation}
Write $p_0=r^--l$.  The wide action set $\mathcal A_{b,r}$ consists of
lists $\alpha=(s_1,\ldots,s_M)$ of valid child grid states satisfying
\[
 |s_i^--l_{i-1}-p_0|\le10b,
 \qquad |s_i^+-l_i-p_0|\le10b.
\]
Consecutive target intervals need not meet.  Their discrepancy is paid
for by the connection charge $\conn_r(\alpha)$ defined in
Section~\ref{sec:geometry}.  For child tables $f$, put
\begin{align}
 V_r(\alpha;f)&=\conn_r(\alpha)+\sum_{i=1}^M f(I_i,s_i),
 \label{eq:action-cost}\\
 \mathcal B_I(f)(q)
 &=\min\left\{w_0(I,q),\;
 \min_{(b,r)}\left[L|q-r|_1+
        \min_{\alpha\in\mathcal A_{b,r}}V_r(\alpha;f)\right]\right\}.
 \label{eq:bellman}
\end{align}
Every minimum over an empty set is $+\infty$.

The operator $\mathcal B_I$ is monotone in its child tables and is
$L$-Lipschitz in $q$.  It also preserves the lower bound
$\max\{c/a,|z|\}$.  For the latter assertion, if each child satisfies
the lower bound, then
\[
 V_r(\alpha;f)
 \ge\frac{\conn_r(\alpha)+\sum_i c(I_i,s_i)}a
 \ge c(I,r)/a,
\]
while the net-change-of-shift inequality gives
$V_r(\alpha;f)\ge|z(I,r)|$.  Endpoint Lipschitz continuity extends
both bounds along the cone, and $w_0$ satisfies them too.

Define the warmup tables at nonleaves by
\begin{equation}\label{eq:warmup}
 w_j(I,q)=\mathcal B_I(w_{j-1})(q)\quad(1\le j\le T),
 \qquad
 m=\frac1T\sum_{j=0}^{T-1}w_j,\qquad d_0=w_T.
\end{equation}
Here the operator receives the indicated tables on children.  Because
$w_1\le w_0$ and the operator is monotone, the warmup is nonincreasing.
In particular,
\begin{equation}\label{eq:warmup-invariants}
 \max\{c/a,|z|\}\le d_0\le m\le w_0\le2Fc,
 \qquad \mathcal B_I(d_0)\le d_0.
\end{equation}
The last inequality follows by one further application of the same
monotonicity.  The tables $w_j$ and $m$ are $L$-Lipschitz.

The average $m$ is useful because shifting its iteration range by one
changes it by only the initial-to-final drop divided by $T$:
\[
 \frac1T\sum_{j=1}^T w_j
 =m-\frac{w_0-w_T}{T}.
\]
For any fixed representative and wide action, average the inequalities
$w_j(I,q)\le L|q-r|_1+V_r(\alpha;w_{j-1})$ and then minimize over
that action. The result is
\[
 m(I,q)-\frac{w_0(I,q)-w_T(I,q)}T
 \le L|q-r|_1+\min_{\alpha\in\mathcal A_{b,r}}V_r(\alpha;m).
\]
Thus parent and child cost proxies satisfy the Bellman comparison up to
a controlled averaging error. Section~\ref{sec:comparison} will sum this
error over a tree of comparison states.

\subsection{Local evaluation of later cones}

The Bellman minimum, and later minima containing selected action cones,
admit the same local evaluation rule.  This rule is also what bounds the
computation when a seed is wrong. We enumerate complete cells because
Section~\ref{sec:bands} will make one shared choice for every cell; that
choice must be independent of which query first requests the cell.

\begin{lemma}[Cone locality]\label{lem:cone-locality}
Fix a good seed as in Theorem~\ref{thm:refinement}.  Take a minimum of
$w_0(I,q)$ and any family of cones of slope $L$ with nonnegative base
costs at representatives $(b,r)$.  If $w_0(I,q)=0$, its value is zero.
Otherwise it suffices to enumerate scales
\begin{equation}\label{eq:local-scales}
 w_0(I,q)/(16F)\le b\le8a w_0(I,q)
\end{equation}
and the following whole cells at each scale.  Partition endpoint-pair
space into half-open cells of side $\theta b$ on each axis, with lower
boundaries at multiples of $\theta b$.  Enumerate the cells intersecting
the closed sup square of radius $w_0(I,q)$ about $q$, and all valid
grid representatives in these cells.
\end{lemma}

\begin{proof}
Write $w=w_0(I,q)>0$.  A competitive cone has value at most $w$ and
therefore $|q-r|_1\le w/L$.  Since $L_0/L=1/4$,
$3w/4\le w_0(I,r)\le5w/4$.  Representative eligibility yields
\[
 b\ge3w/(32F)>w/(16F),\qquad b\le5aw<8aw.
\]
The representative belongs to the enumerated square and hence to an
enumerated cell.  Conversely, enumerating a whole cell only adds cones
already in the global family.  This proves exactness.
\end{proof}

For arbitrary local $w_0$, the scale count in~\eqref{eq:local-scales}
is $O(1+\log(aF))$, polynomial in $\ell$.  There are
$O((1+F/\theta)^2)$ cells per scale and $O((1+P)^2)$ grid states per
cell.  Indeed $w_0(I,q)/b\le16F$, and
$\theta b/g_b\le P$.  Cells and grid points are generated by ranges
of integer indices, without scanning unrelated positions.  These count
bounds do not assume a correct seed.  Empty lists give the $w_0$ option.

We have now obtained nonincreasing propagated estimates and the cost
proxy $m$.  The next step uses $m$ to identify narrow sets of shared
child alignments.  Expensive searches over the wide actions will use
earlier tables; the narrow sets limit calls to current tables.
